\chapter{Experiments \& Evaluation}
Our goal in experimentation was two-fold. It was our intuition that it would be useful to visualize the filter response when applied to acquired \tdd{} with the result of the MSII applied, both regular and irregular synthetic \tdd{} with the Dirac delta function applied.

A smoothing filter is an application of the Heat equation \todoCitation{Jaehne}
Gaussian distribution
Gaussian blurring


Two purposes:
	1. see the behavior of the filter in different data sets, and to
	1. compare the speedup obtained on given different problem sizes and processor counts

Two kinds of data sets:
	synthetic
		square tessallations 2 and 4 triangles
		hexagonal tessalations
		circle triangulation
	acquired


%
%
%
%
%
%
\section{Synthetic Data}
\dots
Flat surfaces with Dirac Delta function
An important test for any filter\todoCitation{} , is to apply the filter iteratively upon a Dirac delta function\todoCitation{}.

%
%
%
%
\subsection{Square Tesselations, Two Triangles}
%\begin{figure}[ht]
	\begin{subfigure}[b]{0.32\linewidth}
		\includegraphics[width=\linewidth]
		{data/synthetic_meshes/square_tesselation_2tri_Dirac_delta_1_v9_f8_wireframe.png}
		\caption{R=1, wireframe}\label{fig:sq2.a}
	\end{subfigure}
	\begin{subfigure}[b]{0.32\linewidth}
		\includegraphics[width=\linewidth]
		{data/synthetic_meshes/square_tesselation_2tri_Dirac_delta_1_v9_f8_funcvals_0iter_crop.png}
		\caption{R=1, convolutions 0}\label{fig:sq2.b}
	\end{subfigure}
	\begin{subfigure}[b]{0.32\linewidth}
		\includegraphics[width=\linewidth]
		{data/synthetic_meshes/square_tesselation_2tri_Dirac_delta_1_v9_f8_funcvals_1iter_crop.png}
		\caption{R=1, convolutions 1}\label{fig:sq2.c}
	\end{subfigure}

	\bigskip
	\begin{subfigure}[b]{0.32\linewidth}
		\includegraphics[width=\linewidth]
		{data/synthetic_meshes/square_tesselation_2tri_Dirac_delta_10_v441_f800_wireframe.png}
		\caption{R=10, wireframe}\label{fig:sq2.d}
	\end{subfigure}
	\begin{subfigure}[b]{0.32\linewidth}
		\includegraphics[width=\linewidth]
		{data/synthetic_meshes/square_tessellation_2tri_Dirac_delta_10_v441_f800_funcvals_0iter.png}
		\caption{R=10, convolutions 0}\label{fig:sq2.e}
	\end{subfigure}
	\begin{subfigure}[b]{0.32\linewidth}
		\includegraphics[width=\linewidth]
		{data/synthetic_meshes/square_tessellation_2tri_Dirac_delta_10_v441_f800_funcvals_100iter.png}
		\caption{R=10, convolutions 100}\label{fig:sq2.f}
	\end{subfigure}
	\caption[Six views, comparing two differently sized of bisected square tessellations]{Comparison of two differently sized bisected square tessellations, generated with parameters R set to 1 and 10: (a) R=1 in wireframe (b) R=1 colored by function value before convolving the filter (c) R=1 colored by function value after convolving the filter once (d) R=10 in wireframe (e) R=10 colored by function value before convoving the filter (f) R=10 colored by function value after convolving the filter 100 times.}
	\label{fig:sq2}
\end{figure}

\todoResearch{Why does sq2 10 need 100 iters to match sq 1 at 1 iters?}
\todoStyle{Why does sq2 1 iters 0 get differnt spacing?}
\todoStyle{Adjusting height, shrinks width => images no longer centered in their columns}

%
%
%
%
\subsection{Square Tesselations, Four Triangles}

%
%
%
%
\subsection{Hexagonal Tesselation}
\begin{figure}[ht]
\ffigbox
	{\begin{subfigure}[b]{0.48\linewidth}
		\includegraphics[width=1.0\linewidth,height=0.3\textheight,keepaspectratio]{data/synthetic_meshes/hexagonal_tessellation_Dirac_delta_1_v31_f42_wireframe.png}
		\caption{Hex v31\_f42 wireframe}\label{fig:hex.a}
	\end{subfigure}
	\begin{subfigure}[b]{0.48\linewidth}
		\includegraphics[width=1.0\linewidth,height=0.3\textheight,keepaspectratio]{data/synthetic_meshes/hexagonal_tessellation_Dirac_delta_10_v1057_f1986_wireframe.png}
		\caption{Hex v1057\_f1986 wireframe}\label{fig:hex.b}
	\end{subfigure}

	\bigskip
	\begin{subfigure}[b]{0.48\linewidth}
		\includegraphics[width=1.0\linewidth,height=0.3\textheight,keepaspectratio]{data/synthetic_meshes/hexagonal_tessellation_Dirac_delta_1_v31_f42_funcvals_0iter_crop.png}
		\caption{Hex v31\_f42 iter 0}\label{fig:hex.c}
	\end{subfigure}
	\begin{subfigure}[b]{0.48\linewidth}
		\includegraphics[width=1.0\linewidth,height=0.3\textheight,keepaspectratio]{data/synthetic_meshes/hexagonal_tessellation_Dirac_delta_10_v1057_f1986_funcvals_0iter_crop.png}
		\caption{Hex v1057\_f1986 iter 0}\label{fig:hex.d}
	\end{subfigure}

	\bigskip
	\begin{subfigure}[b]{0.48\linewidth}
		\includegraphics[width=1.0\linewidth,height=0.3\textheight,keepaspectratio]{data/synthetic_meshes/hexagonal_tessellation_Dirac_delta_1_v31_f42_funcvals_2iter_crop.png}
		\caption{Hex v31\_f42 iter 2}\label{fig:hex.e}
	\end{subfigure}
	\begin{subfigure}[b]{0.48\linewidth}
		\includegraphics[width=1.0\linewidth,height=0.3\textheight,keepaspectratio]{data/synthetic_meshes/hexagonal_tessellation_Dirac_delta_10_v1057_f1986_funcvals_200iter_crop.png}
		\caption{Hex v1057\_f1986 iter 200}\label{fig:hex.f}
	\end{subfigure}}
	{\caption[Synthetic Hexagonal Tessellations, Dirac delta function]{A synthetic hexagonal tessellation, subdivided by triangles, with a Dirac delta function applied: (a) v31 f42 wireframe (b) v1057 f1986 wireframe (c) v31 f42 colored by function value before filter (d) v1057 f1986 colored by function value before filter (e) v31 f42 colored by function value after 2 iterations (e) v1057 f1986 colored by function value after 200 iterations.
% All using the colorramp "Hot (improved)"~\cite[p.~???]{Brewer2003}~\cite[p.~19]{Giga17}, visualized using GigaMesh~\cite{Mara10}, exported as png after disabling the background grid [f7], maximizing the window, disabling screenshot cropping, as well as rejecting tiled rendering, finally cropping to content in GIMP.
	}\label{fig:hex}}
\end{figure}
\todoCitation{}


%
%
%
%
\subsection{Random Circle Trigangulation}
%\begin{figure}[ht]
\ffigbox
	{\begin{subfigure}[b]{0.48\linewidth}
		\includegraphics[width=1.0\linewidth,height=0.3\textheight,keepaspectratio]{data/synthetic_meshes/random_circle_tessellation_Dirac_delta_1_v11_f12_wireframe.png}
		\caption{R.Circ v11\_f12 wireframe}\label{fig:rcirc.a}
	\end{subfigure}
	\begin{subfigure}[b]{0.48\linewidth}
		\includegraphics[width=1.0\linewidth,height=0.3\textheight,keepaspectratio]{data/synthetic_meshes/random_circle_tessellation_Dirac_delta_10_v641_f1252_wireframe.png}
		\caption{R.Circ v641\_f1252 wireframe}\label{fig:rcirc.b}
	\end{subfigure}

	\bigskip
	\begin{subfigure}[b]{0.48\linewidth}
		\includegraphics[width=1.0\linewidth,height=0.3\textheight,keepaspectratio]{data/synthetic_meshes/random_circle_tessellation_Dirac_delta_1_v11_f12_funcvals_0iter.png}
		\caption{R.Circ v11\_f12 iter 0}\label{fig:rcirc.c}
	\end{subfigure}
	\begin{subfigure}[b]{0.48\linewidth}
		\includegraphics[width=1.0\linewidth,height=0.3\textheight,keepaspectratio]{data/synthetic_meshes/random_circle_tessellation_Dirac_delta_10_v641_f1252_funcvals_0iter.png}
		\caption{R.Circ v641\_f1252 iter 0}\label{fig:rcirc.d}
	\end{subfigure}

	\bigskip
	\begin{subfigure}[b]{0.48\linewidth}
		\includegraphics[width=1.0\linewidth,height=0.3\textheight,keepaspectratio,height=0.3\textheight,keepaspectratio]{data/synthetic_meshes/random_circle_tessellation_Dirac_delta_1_v11_f12_funcvals_0iter.png}
		\caption{R.Circ v11\_f12 iter 2}\label{fig:rcirc.e}
	\end{subfigure}
	\begin{subfigure}[b]{0.48\linewidth}
		\includegraphics[width=1.0\linewidth,height=0.3\textheight,keepaspectratio,height=0.3\textheight,keepaspectratio]{data/synthetic_meshes/random_circle_tessellation_Dirac_delta_10_v641_f1252_funcvals_10000iter.png}
		\caption{R.Circ v641\_f1252 iter 10,000}\label{fig:rcirc.f}
	\end{subfigure}}
	{\caption[Synthetic random vertices equally distributed per radius, Dirac delta function]{A synthetic circle filled with random vertices equaly distributed per radius, triangulated by Delauney method~\cite[p.~??]{todoCitation}, with a Dirac delta function applied: (a) v11\_f12 wireframe (b) v641\_f1252 wireframe (c) v11\_f12 colored by function value before filter (d) v641\_f1252 colored by function value before filter (e) v11\_f12 colored by function value after 2 iterations (f) v641\_f1252 colored by function value after 10,000 iterations.
%All using the colorramp "Hot (improved)"~\cite[p.~???]{Brewer2003}~\cite[p.~19]{Giga17}, visualized using GigaMesh~\cite{Mara10}, exported as png after disabling the background grid [f7], maximizing the window, disabling screenshot cropping, as well as rejecting tiled rendering, finally cropping to content in GIMP.
}\label{fig:rcirc}}
\end{figure}
\todoCitation{}
\todoResearch{Why and who equally distributed}



%
%
%
%
%
%
\section{Acquired Data}
\dots

%
%
%
%
\subsection{University Seal}
Unisiegel\_\- UAH\_\- Ebay-Siegel\_\- Uniarchiv\_\- HE2066-60\_\- 010614\_\- partial\_\- ASCII.ply
%\begin{figure}[ht]
	\begin{subfigure}[b]{0.32\linewidth}
		\includegraphics[width=\linewidth]{data/acquired_meshes/unisiegel_wireframe.png}
		\caption{wireframe}\label{fig:unisiegel.a}
	\end{subfigure}
	\begin{subfigure}[b]{0.32\linewidth}
		\includegraphics[width=\linewidth]{data/acquired_meshes/unisiegel_0iter.png}
		\caption{$c=0$}\label{fig:unisiegel.b}
	\end{subfigure}
	\begin{subfigure}[b]{0.32\linewidth}
		\includegraphics[width=\linewidth]{data/acquired_meshes/unisiegel_100iter.png}
		\caption{$c=100$}\label{fig:unisiegel.c}
	\end{subfigure}
	\caption[Three views of the Univeristy of Heidelberg Seal]{Three views of the Univeristy of Heidelberg seal (a) in wireframe, (b) colored by \gls{tMSIIf} value before convolving the filter, (c) colored by function value after convolving the filter 100 times.}
	\label{fig:unisiegel}
\end{figure}


%
%
%
%
\subsection{A Flat surface}
Flat surfaces have NOISE!
Figure~\ref{fig:ILATO}: ILATO\_1A\_SM2066-HE5-60\_070214\_merged\_GMO\_r1.00\_n4\_v256
%\begin{figure}[ht]
	\begin{subfigure}[b]{0.49\linewidth}
		\includegraphics[width=\linewidth]
		{data/acquired_meshes/ILATO_1A_SM2066-HE5-60_070214_merged_GMO_r1_n4_v256_wireframe.png}
		\caption{wireframe}\label{ILATO:bun.a}
	\end{subfigure}
	\begin{subfigure}[b]{0.49\linewidth}
		\includegraphics[width=\linewidth]
		{data/acquired_meshes/ILATO_1A_SM2066-HE5-60_070214_merged_GMO_r1_n4_v256_funcvals_0iter.png}
		\caption{$c=0$}\label{fig:buILATOn.b}
	\end{subfigure}
	
	\bigskip
	\begin{subfigure}[b]{0.49\linewidth}
		\includegraphics[width=\linewidth]
		{data/acquired_meshes/ILATO_1A_SM2066-HE5-60_070214_merged_GMO_r1_n4_v256_funcvals_1000iter.png}
		\caption{$c=1000$}\label{fig:ILATO.c}
	\end{subfigure}
	\begin{subfigure}[b]{0.49\linewidth}
		\includegraphics[width=\linewidth]
		{data/acquired_meshes/ILATO_1A_SM2066-HE5-60_070214_merged_GMO_r1_n4_v256_funcvals_3000iter.png}
		\caption{$c=3000$}\label{fig:ILATO.d}
	\end{subfigure}
	\caption[Four views of a flat surface]{Four views of a flat surface (a) in wireframe (b) colored by MSII function value before convolving the filter (c) colored by function value after convolving the filter 1,000 times (c) and after 3,000 times.}
	\label{fig:ILATO}
\end{figure}


%
%
%
%
\subsection{Stanford Bunny}
Figure~\ref{fig:bun}: http://graphics.stanford.edu/data/3Dscanrep/ (Stanford Bunny)
%\begin{figure}[ht]
	\begin{subfigure}[b]{0.32\linewidth}
		\includegraphics[width=\linewidth]
		{data/acquired_meshes/bun_zipper_edited_r1_n4_v256_wireframe.png}
		\caption{wireframe}\label{fig:bun.a}
	\end{subfigure}
	\begin{subfigure}[b]{0.32\linewidth}
		\includegraphics[width=\linewidth]
		{data/acquired_meshes/bun_zipper_edited_r1_n4_v256_funcvals_0iter.png}
		\caption{$c=0$}\label{fig:bun.b}
	\end{subfigure}
	\begin{subfigure}[b]{0.32\linewidth}
		\includegraphics[width=\linewidth]
		{data/acquired_meshes/bun_zipper_edited_r1_n4_v256_funcvals_100iter.png}
		\caption{$c=100$}\label{fig:bun.c}
	\end{subfigure}
	\caption[Three views of the Stanford Bunny]{Three views of the Stanford Bunny (a) in wireframe (b) colored by function value before convolving the filter (c) colored by function value after convolving the filter 100 times.}
	\label{fig:bun}
\end{figure}


%
%
%
%
%
%
\section{Evaluation}
\ldots

%
%
%
%
\subsection{Compute Times}

Figure~\ref{fig:computeTimesLP} shows how compute times increase linearly with both mesh size and number of iterations in a very predictable way when total compute time is at least 0.1 seconds, and less predictable for shorter periods do to the nature of thread optimization at the processor level and variable memory read times.\todoResearch{add formula for timing (or at least ref to it) here.}\todoResearch{process time noise at very fast speeds}
\begin{figure}[ht]
	\centering
	\includegraphics[width=1.0\linewidth,height=1.0\textheight,keepaspectratio]
		{figures/computeTimesLinespoints.png}
	\RawCaption{\caption[Compute Times - Linespoints]{Compute Times of Applying
		the	One-Ring Filter for Selected Numbers of Iterations onto Acquired and
		Synthetic 3D Meshes of Varying Sizes}
		\label{fig:computeTimesLP}}
\end{figure}

Figure~\ref{fig:computeTimesS} shows how compute times increase with both meshsize and number of iterations.\todoCitation{wikipedia Euler characteristic Polyhedra}
\begin{figure}[ht]
	\centering
	\includegraphics[width=1.0\linewidth,height=1.0\textheight,keepaspectratio]{figures/computeTimesScatter.png}
	\RawCaption{\caption[Compute Times - Scatter]{Compute Times for Different
		Hardware Configurations by increaseing Mesh Size and Filter Iterations}
		\label{fig:computeTimesS}}
\end{figure}

%
%
%
%
\subsection{Curious Trend, ratio goes to 2}

\begin{figure}[ht]
	\centering
	\includegraphics[width=1.0\linewidth,height=1.0\textheight,keepaspectratio]{figures/numFacesByVerticesGoTo2.png}
	\RawCaption{\caption[Ratio of Faces / Vertices]{Ratio of Faces to Vertices by
		Increasing Vertex Count}
		\label{fig:ratioFacesVertices}}
\end{figure}



%\subsection{Debossed H}
%In Figure \ref{fig:h}, we show a debossed capital letter H.\footnote{The H is
%as a nod to Heidelberg University and the cuniform script studied by the FCGL.}
%\begin{figure}[ht]
%\centering
%	\begin{subfigure}{.48\linewidth}
%		\centering
%		\resizebox{0.48\linewidth}{!}{\input{data/synthetic_meshes/h.tikz_labels.tex}}
%		\caption{HWireframe}\label{fig:h.a}
%	\end{subfigure}
%	\hfill
%	\begin{subfigure}{.48\linewidth}
%		\centering
%		\includegraphics[width=0.48\linewidth]{data/synthetic_meshes/h_colored.png}
%		\caption{HColored}\label{fig:h.b}
%	\end{subfigure}
%	\caption[A debossed H, which contains 22 vertices and 36 faces.]{A debossed H,
%	which contains 22 vertices and 36 faces: (a) wireframe (b) colored by the
%	relation to its distance to an underlying plane, in RdGy
%	colorramp~\cite[p.~???]{Brewer2003}~\cite[p.~19]{Giga17}, visualized using the
%	GigaMesh~\cite{Mara10} framework with triangle edges rendered.}\label{fig:h}
%\end{figure}
%To evaluate methods available for discrete surfaces, we can increase the number of vertices of our synthetic wedge using five iterations of the mid-edge subdivision scheme [PR97,HW99].~\cite[p.~38]{Mara12}

%\subsection{Mars Crater}
%Mars dataset crater as a Digital Terrain Models (DTMs) Mention in Mara 3.6
%Summary “Dali” inspired methodProcessing regular grids like Digital Terrain Models (DTMs) will gain dramatic performance increases using the estimator, while processing irregular grids with high curvatures will strongly benefit from precise computation of the volume integral invariant.~\cite[p.~143]{Mara12}

%
%
%
%
%
%
\section{Summary}
\dots
